\section[Conclusão]{Conclusão}

% Neste item o aluno deverá refletir sobre:
% \begin{itemize}
%     \item Crítica e autocrítica em relação a sua atuação no projeto na parte técnica como na parte de soft skills.
%     \item Comente o relacionamento entre as disciplinas cursadas e a AGES.
%     \item Este projeto foi o melhor projeto trabalhado? Justifique.
%     \item Relate as lições aprendidas (retrospectiva pessoal):
%           \begin{itemize}
%               \item O que foi positivo no projeto.
%               \item O que podia ter sido melhor (em ternos de banco, desenvolvimento, arquitetura, soft skills....)
%           \end{itemize}
% \end{itemize}
% (No mínimo uma página de relato)

Refletindo sobre minha experiência com os colegas e as tecnologias utilizadas neste projeto, considero que esta AGES 
foi uma oportunidade de aprendizado e colaboração extremamente valiosa. Desde o início, me senti motivado a ser um colega 
prestativo e a oferecer suporte àqueles que estavam enfrentando os mesmos desafios que eu já havia enfrentado no passado.

O time conseguiu entregar todas as \textit{features} prometidas à cliente, exceto por uma pequena task de integração, graças 
ao apoio dos AGES IV e do nosso único AGES III. No entanto, apesar desses pontos positivos, considero que muitas dificuldades 
e empecilhos do projeto surgiram de decisões relacionadas à gestão do trabalho.

Nas sprints 0 e 1, fomos divididos em squads para diferentes áreas de desenvolvimento, o que desagradou alguns iniciantes que foram 
alocados em áreas onde não queriam atuar. Esses integrantes demonstraram desmotivação ao longo do semestre e, somado a fatores comuns 
entre AGES I, como timidez e inexperiência, acabaram contribuindo menos do que o esperado. 

Sinto que, com mais \textit{feedbacks} tanto da minha parte quanto dos insatisfeitos, poderíamos ter organizado melhor a ideia de grupos dentro do 
time, resultando em uma joranda mais positiva no geral. Digo isso porque, a partir da sprint 2, as squads foram desfeitas, e cada aluno 
passou a trabalhar individualmente em tasks que conseguisse assumir. Uma das minhas maiores críticas à minha experiência na AGES I é, 
justamente, esse modelo de trabalho. 

Apesar das críticas levantadas, não vejo este projeto de forma negativa, muito pelo contrário. Minhas observações refletem minha preocupação 
em tentar realizar um bom trabalho, especialmente porque me diverti aprendendo novas tecnologias e colaborando com colegas e uma professora orientadora 
excepcionais.

Quanto às lições aprendidas, elas se assemelham às da minha experiência na AGES I, mas com um tom diferente. Na primeira passagem, destaquei a timidez e 
as dificuldades de comunicação como \textit{"um dos aspectos de maior relevância do trabalho em grupo"}. Desta vez, sinto que me comuniquei bastante, mas ainda assim 
poderia ter feito mais. Os \textit{feedbacks} levantados nas retrospectivas das sprints precisam ser levados mais a sério para resolver qualquer empecilho assim que surgir. 
Além disso, percebi que me acostumei a trabalhar com colegas específicos, possivelmente deixando outros sobrecarregados sem perceber. Reitero, portanto, que a comunicação 
frequente continua sendo, na minha visão, o aspecto mais importante do trabalho em equipe.  
